\subsection{缠绕算子}

接下来，我们想解释：考察图 \ref{f:table} 中出现的范畴之间的缠绕算子，如何重新得到若干基本定理，并产生许多引人注目的猜想。

在这一点上，前面许诺的对局部几何朗兰兹对应朴素表述 \eqref{e:loclangnaive} 的技术性修正，自然就以如下方式显现出来。

\ssszh 回忆 $p$-adic 情形中最简单的表示是非分歧表示。令人高兴的是，在非分歧情形中，我们已能遇到几何猜想朴素表述 \eqref{e:loclangnaive} 所需的修正。

为看清这一点，回忆局部朗兰兹应当把仿射 Grassmann 簇上的 D-模 $\Dmod(\Gr_G)$，与平凡丛上的拟凝聚层 $\QCoh(\on{pt}/\cG)$ 互换。特别地，二者的自同态应当典范地等同；我们来看看它们是什么。

在自守侧，回忆 $\Gr_G$ 是陪集空间 $LG/L^+G$，如第 \ref{s:catrepgrp} 节所述，有
\begin{equation} \label{e:dersat1}
 \Hom_{\Dmod(LG)\mod}( \Dmod(\Gr_G), \Dmod(\Gr_G)) \simeq \Dmod(L^+G \backslash
 LG / L^+G),
\end{equation}
即双球层的导出 Satake 范畴。作为对初次接触此事的读者的可能提醒：这个范畴不是相应阿贝尔范畴的导出范畴，而更为丰富。

在谱侧，利用映射 $\cfg / \cG \rightarrow \LS$ 在平凡丛处形式 étale 这一事实，可以如例 \ref{e:fibprodvar} 那样算出
\begin{equation} \label{e:dersat2}
\begin{split}
&\Hom_{\TwoCoh(\LS)}(\QCoh(\on{pt}/\cG), \QCoh(\on{pt}/\cG)) \\
&\qquad \simeq \QCoh( \on{pt}/\cG \underset{\LS} \times \on{pt}/\cG)
\simeq \QCoh( \on{pt} \underset{\cfg} \times \on{pt} / \cG).
\end{split}
\end{equation}
这里需要两点解释。第一，对 \eqref{e:dersat2} 中最后的等同，可以用下面这条一般观察。设有代数群 $K$ 作用在簇 $X, Y, Z$ 上。给定 $K$-等变的对应
$$X \rightarrow Z \leftarrow Y,$$
则有商（导出）叠的典范同构
\[
      X/H \underset{Z/H} \times Y/H \simeq (X \underset Z \times Y)/H.
\]
今后我们会反复使用这一事实，不再另行说明。

第二，提醒读者：\eqref{e:dersat2} 中出现的纤维积是导出概形。具体地说，这只是意味着我们同时记住 Tor；于是自交
$$\on{pt} \underset{ \fg} \times
\on{pt}$$
是用 Koszul 消解作标准计算得到的、外代数 $\Sym \cfg^*[1]$ 的谱。

因此，我们可以朴素地期望 \eqref{e:dersat1} 与 \eqref{e:dersat2} 之间有幺半范畴的等价。其实，作为 Bezrukavnikov--Finkelberg \cite{bezfinksat} 的一个重要定理，导出 Satake 等价把 \eqref{e:dersat1} 与 \eqref{e:dersat2} 的一个放大等同起来。更确切地说，他们证明了等价
\begin{equation} \label{e:dersat} \Dmod(L^+G \backslash LG / L^+G) \simeq
\ICn( \on{pt} / \cG  \underset{
\cfg / \cG} \times \on{pt} / \cG), \end{equation}
其中右端是带{\em 幂幺奇异支撑}的归纳凝聚层构成的幺半范畴。

\ssszh 我们简要说明什么是归纳凝聚层，以及施加奇异支撑条件是什么意思。它们算是对拟凝聚层这一想法的一些技术性修正，初次接触这一圈想法的读者可以放心跳过。我们先讲归纳凝聚层，再转向奇异支撑条件。

用最朴素的话说，归纳凝聚层的理论源自有限生成模与有限生成投射模（即向量丛）之间的差别。在阿贝尔范畴中，比如对仿射概形 $X$，向量丛当然比有限生成模少，除非 $X$ 是若干约化点的无交并。

然而在导出范畴中，这一区别更为微妙。事实上，若 $X$ 光滑，则任何有限生成模都容许由向量丛给出的有限消解。因此，人们无法区分有限生成模的有限复形与向量丛的有限复形。若 $X$ 不光滑，它们在导出范畴中又有差别：一般的有限生成模此时需要无限投射消解。具体地说，可以取 $X$ 的任一奇点处的摩天层。特别地，这两类复形之间的差别给出了 $X$ 奇点程度的一个度量。

\ssszh 我们按照 Krause \cite{krause} 与 Gaitsgory \cite{gaitsgoryindcoh} 的做法，把前面几段讲得更仔细些。首先回忆：对代数簇 $X$，其拟凝聚层范畴 $\QCoh(X)$ 由完美复形紧生成，即由局部等价于向量丛的有限复形的对象紧生成。若 $X$ 仿射，这就化归为同义反复：$O_X\mod$ 由正则函数环 $O_X$ 生成。对非仿射的 $X$，所断言的在拟射影簇情形容易，但一般而言是 Thomason \cite{thomasontrobaugh} 一个颇为精细的定理。

作为有限生成模的有限复形的整体类比，考虑凝聚层的有界导出范畴
$$\on{Coh}(X)^b \subset \QCoh(X).$$ 按定义，$\on{IndCoh}(X)$ 是它的归纳完备化，即 $\on{IndCoh}(X)$ 由 $\on{Coh}(X)^b$ 紧生成。显然，$\on{IndCoh}(X)$ 的对象是上同调凝聚的有界复形的形式余极限，而同态复形按复形的同伦极限计算
\[
\Hom_{\on{IndCoh}(X)}( \varinjlim_i \sE_i, \varinjlim_j \sF_j) \simeq
\varprojlim_i
\varinjlim_j \Hom_{\on{QCoh}(X)}( \sE_i, \sF_j).
\]
对任何凝聚层 $\sE$，此式迫使 $\Hom_{\on{IndCoh}(X)}(\sE, -)$ 与余极限交换；而在 $\QCoh(X)$ 中，除非 $\sE$ 完美，这一点不成立。

由于完美复形有界凝聚，把它们的包含函子作归纳完备化，就得到伴随对
\[
\on{QCoh}(X) \rightleftarrows \on{IndCoh}(X),
\]
其中左伴随全忠实，且当 $X$ 光滑时是等价。特别地，人们把 $\on{IndCoh}(X)$ 看作 $\on{QCoh}(X)$ 的放大，并且任一归纳凝聚层都有到拟凝聚层的截断。此外，$\on{IndCoh}(X)$ 带有自然的 $t$-结构，并在有下界的对象上与 $\on{QCoh}(X)$ 一致。不正式地说，$\on{IndCoh}(X)$ 是通过添加上同调次数为 $-\infty$ 的对象得到的。

\begin{re} 回忆一下，我们在讨论 Kac--Moody 表示与环路群上的 D-模时，已经遇到过这种对无限上同调维数的导出范畴的重正规化，参见例 \ref{ex:kmmods}；这或许有助于定位。\end{re}

\begin{re} 在奇点簇上系统地区分完美复形与有界凝聚复形这一想法，此前也出现在代数几何中。例如，商 $\on{IndCoh}(X)/\QCoh(X)$ 就是 Orlov 奇点范畴 \cite{orlovsingularity} 的归纳完备化。\end{re}

\ssszh 接下来解释奇异支撑的想法。它由 Gulliksen \cite{gulliksen}、Eisenbud \cite{eisenbud}、Avramov 与 Buchweitz \cite{avramovbuchweitz} 针对完全交发展起来，又由 Arinkin--Gaitsgory 推广到导出完全交 \cite{arinkingaitssingsupp}；另可参见重要的 \cite{bensoniyengarkrause}。按定义，归纳凝聚层与拟凝聚层之间的差别源自 $X$ 的奇点；特别地，商 $$\IC(X)/\QCoh(X)$$ 将沿 $X$ 的奇异轨迹支撑。给定凝聚层 $\sF$，人们可能想要一个不变量，粗略地记录它在奇点处的何处不再是完美的。奇异支撑对这一问题的回答更为精细。

下面是稍微精确一点的表述。设 $X$ 是拟光滑导出概形，即其余切复形 $T^*(X)$ 支撑在上同调次数 $-1$ 与 $0$ 处。则对上述对象 $\sF$，其奇异支撑是 $H^{-1} T^*(X)$ 的一个闭锥形子集 $\Lambda$。奇异支撑逐仿射地度量 $\OO_X$ 的元素的局部幂零性，从而恢复通常的支撑，以及 $\sF$ 的某些自然的二次自映射。此外，对任何这样的子集 $\Lambda$，都得到全子范畴
\[
   \IC_\Lambda(X) \subset \IC(X)
\]
由奇异支撑含于 $\Lambda$ 的归纳凝聚层构成。例如，若 $\Lambda$ 是零截面，这恢复 $\QCoh(X)$。

虽然我们不打算给出计算奇异支撑的一般定义，但马上会讨论几个特例。不过先给出几条帮助定位的注。

\begin{re}\label{r:singsupp}导出概形拟光滑，当且仅当它局部可写成纤维积 $X \times_Z Y$，其中 $X,Y,Z$ 都是光滑簇。在这个局部图景中，余切复形有如下显式形式。确切地说，对纤维积的一点 $(x,y)$，设其公共像为 $z$，则有切复形的 distinguished 三角形
\begin{equation} \label{e:dts}
    T_{(x,y)} (X \times_Z Y) \rightarrow T_x(X) \oplus T_y(Y) \rightarrow
    T_z(Z) \xrightarrow{+1}.
\end{equation}
对不熟悉这个 distinguished 三角形的读者，先理解 $X$ 与 $Y$ 横截地映入 $Z$ 的情形会有帮助。此时纤维积仍然光滑，上述 distinguished 三角形变成一个（希望是）直观的短正合列。在一般情形，要点在于：纤维积上的导出结构修正了它的切复形，使得横截情形下最朴素的修正依然成立。

把 \eqref{e:dts} 对偶化，可以看出余切丛的 $H^{-1}$ 具体就是那些与 $T_x(X)$、$T_y(Y)$ 的像正交的余方向 $\xi \in T^*_z(Z)$；更不正式地说，就是在 $Z$ 中还剩下可以走的方向。

上述情形的一个重要特例是形如
\begin{equation}
\label{e:vanlocus} \{0\}
\times_{\mathbb{A}^n} Y,\end{equation}
的纤维积，即 $Y$ 上 $n$ 个函数 $f_1, \ldots, f_n$ 的导出零点集。由此可见，拟光滑性是导出代数几何中「隐藏光滑性」哲学的一例。也就是说，纤维积底层的经典概形可以任意奇异，但它所承载的导出概形表现得像完全交，即只是轻度奇异。

\end{re}

下面用几个例子说明奇异支撑究竟是什么。

\begin{ex}\label{ex:singsupp1}对凝聚层 $\sF$，其奇异支撑与零截面的交就是通常的支撑。\end{ex}

\begin{ex}\label{ex:singsupp2}设 $X$ 可实现为光滑簇 $U$ 中的超曲面。此时 $$H^{-1} T^*(X)$$ 由 $X$ 本身（看作零截面）连同放在 $X$ 每个奇点 $x$ 处的一条直线 $\ell_x$ 构成。设凝聚层 $\sF$ 的通常支撑含有这样的 $x$。则其奇异支撑含有 $\ell_x$，当且仅当 $\sF$ 在 $x$ 的任一 Zariski 邻域内都不与完美复形同构。\end{ex}

\begin{ex}设 $X$ 如导出 Satake 的谱侧那样，由向量空间中原点的自交给出：$$X = 0 \underset V \times 0.$$ 此时有 $T^*_{(0,0)} X \simeq V^*[1]$。这里，奇异支撑可以理解为 Koszul 对偶另一侧的通常支撑。确切地说，摩天层 $k$ 生成 $\IC(X)$，取它的自同态就得到等价
\[
    \IC(X) \simeq \Sym V[-2] \mod.
\]
经由这个等同，对任何凝聚层 $\sF$，其奇异支撑与相应 $\Sym V[-2]$-模的通常支撑一致（忽略分次的平移）。这里可以把 $0$ 与 $V$ 换成任意簇的任一光滑点。\end{ex}

\ssszh 讨论了归纳凝聚层与奇异支撑的基础之后，我们回到导出 Satake 等价，即
\begin{equation} \label{e:der} \Dmod(L^+G \backslash LG / L^+G) \simeq
\ICn( \on{pt} / \cG  \underset{
	\cfg / \cG} \times \on{pt} / \cG). \end{equation}
在这个等价中，下标「nilp」指的是归纳凝聚层中奇异支撑落在幂零锥
$$\sN \subset \cfg^* \simeq
H^{-1}T^*( \on{pt}/\cG \underset {\cfg / \cG}   \times \on{pt}/\cG)$$
之中的子范畴。

\ssszh 我们马上会看到，非分歧情形中遇到的模式——必须放大谱侧，允许带幂幺奇异支撑的归纳凝聚层——在另一些匹配缠绕算子的情形中也会出现。这就迫使我们对局部几何朗兰兹的谱侧作调整，办法是从范畴的拟凝聚层过渡到带幂幺奇异支撑的范畴归纳凝聚层
\[
\TwoCoh(\LS) \rightsquigarrow \TwoIC(\LS).
\]
我们马上会谈这指的是什么；目前只需说，有了这一修正，就得到局部几何朗兰兹猜想的现有表述，即等价
\begin{equation}
\label{e:loclang}\Dmod(LG)\mod \simeq \TwoIC(\LS).
\end{equation}

\ssszh 我们简要勾勒带奇异支撑的范畴归纳凝聚层理论，它属于 Arinkin \cite{arinkinindcohcats}。

当范畴层级从归纳凝聚层升到范畴归纳凝聚层时，要建立奇异支撑理论，人们不再要求底簇 $X$ 拟光滑（即余切复形在次数 $-1$ 与 $0$），而要求它光滑（即余切复形在次数 $0$）。原因马上可见。

于是设 $X$ 是光滑代数簇。利用 1-仿射性可以证明：对任何从光滑簇出发的紧合映射 $f: Z \rightarrow X$，$\QCoh(Z)$ 紧生成 $\TwoCoh(X)$ 中支撑于 $\on{im}(f)$ 之上的对象构成的全子范畴。特别地，这个子范畴等同于卷积代数
\[
  \QCoh(Z \underset X \times Z)
\]
上的模。

$X$ 上范畴归纳凝聚层的 $(\infty,2)$-范畴有如下基本性质。第一，它带有 $\TwoCoh(X)$ 的幺半作用，因而可以谈论对象的支撑。此外，它由形如 $\QCoh(Z)$ 的对象（$Z$ 如上）生成。然而，此类对象的自同态现在提升为
\[
   \IC( Z \underset{X}{\times} Z),
\]
粗略地说，随着 $Z \rightarrow X$ 越来越远离光滑映射，它相对于拟凝聚的类似物会变得更大。更确切地说，如注 \ref{r:singsupp} 所述，这个卷积范畴中对象的奇异支撑，沿 $X$ 中与 $Z \times Z$ 的切丛正交的余方向分布。因此，$f$ 在给定点 $x$ 上越不光滑，$\on{2-IndCoh}(X)$ 中由 $\QCoh(Z)$ 生成、且支撑于 $x$ 之上的范畴归纳凝聚层就越多。

特别地，$\QCoh(Z)$ 不再生成 $\on{2-IndCoh}(X)$ 中对应于支撑于 $\on{im}(f)$ 之上的对象的全子范畴，而是生成如下微局部细化后的子范畴。\footnote{确切地说，要看出细化是必要的，注意：若没有细化，取 $Z = X$ 就会把 $\on{2-IndCoh}(X)$ 与 $\TwoCoh(X)$ 等同起来。} 由于假定 $f$ 紧合，可以构造余法簇 $N^\vee_{Z}$，即由沿 $$Z \underset X \times T^*X \rightarrow T^*X$$ 投影 $f^* T^*X \rightarrow T^*Z$ 的核所得到的、$T^*X$ 的闭各向同性子簇。可以给任何范畴归纳凝聚层配上它在 $T^*X$ 中的支撑，而 $\QCoh(Z)$ 生成微支撑落在 $N^\vee_Z$ 之中的对象构成的全子范畴。

更一般地，给定 $T^*X$ 的任一闭锥形各向同性子簇 $\Lambda$，可配上全子范畴
$$
\on{2-IndCoh}_\Lambda(X) \subset \on{2-IndCoh}(X).
$$
它同样由 $\QCoh(Z)$ 的（截断）生成，自同态为 $$\on{IndCoh}_\Lambda(Z \underset X \times Z),$$ 其中 $\Lambda$ 现在用来对通常的归纳凝聚层施加通常的奇异支撑限制，如注 \ref{r:singsupp}。特别地，取 $\Lambda$ 为零截面，就恢复 $\TwoCoh(X)$。

\begin{re} 归纳凝聚层与范畴归纳凝聚层的形式体系还有另一种进路，其中包含分别超出拟光滑与光滑情形的奇异支撑理论；该进路正由 di Fiore--Stefanich 发展中。\end{re}

\ssszh 前面的讨论并不字面地适用于 $\LS$。粗略地说，原因在于它无穷维，不是光滑 Artin 叠。不过人们预期，它带有类似的范畴归纳凝聚层理论。

承认这种形式体系存在，则由 \eqref{e:loclang}，只剩一件事要解释：奇异支撑条件「nilp」指的是什么。不过，鉴于我们在非分歧情形已经遇到的东西，这应当不足为奇。确切地说，$\LS$ 在联络 $(\check{\sP}, \nabla)$ 处的余切复形，由计算余伴随丛
\begin{equation} \label{e:datboi}
   \check{\sP} \overset{\cG} \times \cfg^*
\end{equation}
在穿孔圆盘上平坦截面的 de Rham 复形给出。特别地，余切空间即其零次上同调，典范地等同于 \eqref{e:datboi} 的平坦截面。于是，「nilp」在每个余切空间内切出那些幂幺的平坦截面。
